%============================================================================
%  SKELETON DRAFT — Memory transfer and continual learning in
%  predictive-coding associative networks
%
%  This is a STRUCTURE-ONLY skeleton: section/subsection titles and one-line
%  notes on what goes where. No content is developed. Preamble/skeleton follows
%  the PLOS-style template (Saighi & Rozenberg) and the section ordering of
%  Tang et al. 2023 (Results-first, declarative headings, theorems stated inline
%  in a Model/analysis block, full proofs in Materials and methods).
%
%  Full derivations live in ../maths/detailed_proofs.md .
%  Build:  pdflatex manuscript && bibtex manuscript && pdflatex manuscript && pdflatex manuscript
%============================================================================
\documentclass[10pt,letterpaper]{article}
\usepackage[top=1in,bottom=1in,left=1.1in,right=1.1in]{geometry}

\usepackage{amsmath,amssymb,amsthm}
\usepackage{graphicx}
\usepackage[table]{xcolor}
\usepackage{caption}
\usepackage{subcaption}
\usepackage{cite}
\usepackage[right]{lineno}
\usepackage{fancyhdr}
\usepackage[colorlinks=true,allcolors=blue]{hyperref}

% ---- theorem environments (the paper's formal results) ----
\theoremstyle{plain}
\newtheorem{theorem}{Theorem}
\newtheorem{lemma}{Lemma}
\newtheorem{proposition}{Proposition}
\newtheorem{corollary}{Corollary}
\theoremstyle{remark}
\newtheorem{remark}{Remark}

% ---- captions: bold "Fig" label, PLOS-like ----
\captionsetup{labelfont=bf,labelsep=period,justification=raggedright,singlelinecheck=off}
\renewcommand{\figurename}{Fig}

% ---- header / footer ----
\pagestyle{fancy}
\fancyhf{}
\rfoot{\thepage}
\lfoot{Draft skeleton \today}
\renewcommand{\headrulewidth}{0pt}

% ---- a lightweight figure-placeholder box (no image files needed yet) ----
\newcommand{\figplaceholder}[1]{%
  \begin{center}\setlength{\fboxsep}{14pt}%
  \fbox{\parbox{0.86\textwidth}{\centering\itshape [figure placeholder]\\ #1}}%
  \end{center}}

\setlength{\parindent}{0.5cm}

\begin{document}

%============================================================================
%  TITLE BLOCK
%============================================================================
\begin{flushleft}
{\Large\textbf{Spontaneous memory transfer enables continual learning in adversarially coupled predictive-coding networks}}
\newline\\
Saighi Paul\textsuperscript{1,*}
\\
\bigskip
\textbf{1} Universit\'e Paris-Saclay, CNRS, Laboratoire de Physique des Solides, 91405, Orsay, France\\   
\bigskip
* paul.saighi@universite-paris-saclay.fr
\end{flushleft}

%============================================================================
%  ABSTRACT / AUTHOR SUMMARY  (placeholders only)
%============================================================================
\section*{Abstract}

% don't talk about memory transfer in the whole paper but about the transfer of memory content.

During sleep and quiet wakefulness, the brain spontaneously reactivates neural activity patterns associated with past experience, a phenomenon known as reactivation. Reactivation is considered to support systems consolidation, the process by which the retrieval of memory content becomes progressively more dependent on neocortical rather than hippocampal circuits. This reorganization of memory representations has, in turn, been proposed as a potential mechanism supporting continual learning, the ability to acquire new knowledge over time without forgetting crucial information previously learned.

Reactivation-based systems consolidation has been the subject of extensive modeling efforts, including biologically plausible models in which reactivation arises autonomously from intrinsic network dynamics. 

\textbf{However, these models generally do not derive reactivation priorities from the interaction between the memory systems themselves, leaving reactivation unprioritized or biasing it according to factors such as recency, encoding strength, familiarity, salience, or reward. Here, we show that reactivation priorities can instead emerge from the dialogue between the neural populations involved in consolidation.}

[Then we need to say why we propose this approach. First on a the computational side it is particularly efficient allowin fast near exhaustive memory transfer between subsystem. secondly, when it comes to explaining observations in biology, it is parsimonious, as it allow to explain recency effects (replay of recent events), but also replay of temporally remote but related memory to new learning during system consolidation. Then it is an effective hypotheses about non trivial interaction observation between prefrontal cortex modulation of hippocampal Cornu Ammonis 1 (CA1) activity during sleep. And finaly it is embedded in the predictive coding framework which ] 

%\textbf{Here we show that, given structural modifications motivated by biological observations, the predictive coding (PC) framework provides a natural account of both. (structural modification is weird, find better wording. also PC framework account for both ? not really because we need the structural modification. we need a single elegant and short way to rephrase that that keep its stenght while making it more exact.)} We couple two PC associative memory networks, a teacher holding a given memory content and a student that lacks it, through a shared prediction-error interface. We show that reversing the sign with which the prediction-error interface influences the teacher's activity, which we call adversarial coupling, turns typical pattern-completion dynamics into an autonomous memory-transfer process. The teacher preferentially reactivates components of its stored content that are poorly represented by the student, which progressively acquires them in its own recurrent weights. We show analytically that this reactivation is prioritized according to the student's current representational deficit, so that the process requires neither an explicit replay scheduler nor an external cue.

%Building on this mechanism, we propose a continual-learning architecture in which new memory content is incorporated through local plasticity while interference with previously stored content is mitigated, even when successive memories are highly correlated. The specific way in which teacher and student interact in our model relates to recent observations about prefrontal cortex modulation of hippocampal Cornu Ammonis 1 (CA1) activity during sleep.

\linenumbers

%============================================================================
%  1. INTRODUCTION
%============================================================================
\section{Introduction}

[USING A CAREFULL RIGOROUS SCIENTIFIC TONE, WE NEED VERIFIED CITATIONS IN THIS PART]

A substantial fraction of the brain's metabolic budget is devoted to intrinsic activity \cite{raichle2002budget,raichle2006dark}. Intrinsic activity can be defined as any activity which is not dedicated the processing of current sensory signal or the control of overt behaviour \cite{raichle2002budget,raichle2006dark}. This intrinsic activity is especially evident during sleep and quiet wakefulness, when interaction with the environment is reduced \textbf{yet the brain sustains highly structured spontaneous dynamics (weird flow too unformal)} \cite{steriade1993thalamocortical,diekelmann2010memory}. Morover, The brain’s overall metabolic demand remains substantial throughout sleep: glucose utilization decreases of only one quarter during non-rapid-eye-movement (NREM) sleep but returns towards waking levels during rapid-eye-movement (REM) sleep \cite{buchsbaum1989sleep,maquet1990cerebral}. \textbf{Furthermore (maybe change)}, sleep and sleep-like states are widely conserved across animal lineages, with states satisfying their defining behavioural criteria described in widely separated taxa, including insects and nematodes \cite{cirelli2008essential}. The coexistence of substantial metabolic demand, broad phylogenetic conservation, and the absence of immediate behavioural outcomes has motivated a long-standing effort to identify the functions supported by this activity. 

% correct definitio of intrinsic activity. Improve the formulation of the absence of immediate behavioural consequences.

Among the events taking place during the brain's intrinsic activity, \textbf{one of the best characterised (too unformal)} is neural reactivation: the reinstanciation, during sleep or quiet rest, of activity patterns observed during waking experience. Observations in the rodent hippocampus showed that place cells which was co-active while an animal occupied particular locations during awake activity exhibited an increased tendency to be co-active again during subsequent sleep \cite{wilson1994reactivation}. Replay has since been reported in humans during post-task rest, including tasks with no spatial component \cite{schuck2019sequential,liu2019human}. Subsequent works established that replay is not confined to the hippocampus: reactivation of patterns observed during behavioural tasks has also been shown in the ventral striatum, prefrontal cortex and primary visual cortex \cite{pennartz2004striatum,euston2007pfc,jiwilson2007}. Hippocampal reactivation is commonly observed during sharp-wave ripples, a subset of these events being temporally coordinated with neocortical replay, during slow oscillations and thalamocortical spindles, allowing reactivation in the hippocampus and neocortex to be aligned within a common temporal scaffold \cite{klinzing2019mechanisms,staresina2023coupled}. This coordination is thought to provide windows favourable to memory processing and is commonly described as a hippocampo-cortical dialogue (REF).

This dialogue is often considered as supporting system consolidation, the process by which the retrieval of memory during cognitive tasks becomes progressively less dependent on the hippocampus and more dependant on the neocortical circuits \cite{klinzing2019mechanisms, diekelmann2010memory, squire2015consolidation}. Moreover, this reorganization has been associated to the maturation of memory engram in the prefrontal cortex \cite{kitamura2017engrams}. In this work, we frame system consolidation as a form of memory transfer, whereby memory content is partially or completly transfered from a subsystem to another.
Replay, and more broadly spontaneous activity during sleep, has been assigned many other functions, among which we have : Synaptic renormalisation and downscaling \cite{tononi2014price, tononi2014sleep}, the extraction of statistical regularities and generalisation of learning across episodes \cite{klinzing2019mechanisms}, and metabolite clearance \cite{xie2013sleep}. These views are not mutually exclusive and this work specifically address the memory transfer phenomenon.

the complementary-learning-system (CLS) framework gives memory transfer a computational role for continual learning (CL). Continual learning can be defined as the ability for a system to acquire new knowledge over time while retaining crucial information previously learned. 
In the CLS, the Hippocampus act as a fast learning subsystem, allowing the representations proper to new experiences to be quickly stored without immediately forcing their storage as cortical representations.Then a slow distillation process, 


Multiple hypothesis have been formulated on how this slow distillation preserve the system from catastrophic forgetting. \textbf{An increasingly popular view}, supported by recent observations (REF), is that during replay, the reactivation of new memory traces expressed by the Hippocampus are interleaved with the reactivation of older memory traces expressed by the cortex.


This slow distillation is considered to allow the mitigation of inteferences with old memory representations during the integration of new information.


% The complementary-learning-systems framework gives this transfer a computational role in continual learning. Rapid hippocampal plasticity allows individual experiences to be encoded without immediately forcing them into existing cortical representations, whereas slower cortical plasticity supports the gradual extraction of shared structure \cite{mcclelland1995cls,kumaran2016what}. The separation is important because learning correlated patterns sequentially can cause updates induced by a new pattern to degrade representations acquired earlier, a form of catastrophic interference \cite{mccloskey1989catastrophic,robins1995catastrophic}. Interleaving new information with reactivations of older representations constrains each update by both, thereby reducing destructive crosstalk. Computational studies further suggest that preferentially interleaving older memories that overlap with new content can substantially reduce the amount of rehearsal required \cite{saxena2022interleaved}. We therefore examine the specific hypothesis that interleaved communication among complementary associative memories can integrate new, correlated content while preserving previously stored structure.

%Why memory should be organised in two stages is addressed by complementary learning systems theory \cite{mcclelland1995cls,kumaran2016what}. A single network with distributed representations that learns items in sequence suffers catastrophic interference, the weight changes encoding a new item degrading those that support earlier ones \cite{mccloskey1989catastrophic,ratcliff1990connectionist}. Interleaving new items with old ones during learning removes most of this interference, but requires the old items to remain available while the new ones are learned. The hippocampus is proposed to supply that availability: it encodes experience rapidly, and its offline reactivation interleaves new content with content the neocortex already holds, so that a single cortical substrate can integrate new and potentially correlated memory content without overwriting what it stores \cite{mcclelland1995cls,robins1995catastrophic,mcclelland2020integration}. Other accounts of how sleep protects memory emphasise global synaptic downscaling \cite{tononi2014price} or schema-dependent integration that bypasses slow interleaving \cite{tse2007schemas}. We adopt the interleaved-replay account; the interleaving, and the cancellation of interference it produces, are exhibited directly in the network dynamics below.

%The functional significance usually attached to this transfer is continual learning. Distributed connectionist substrates that acquire new content by gradient-based, error-driven plasticity overwrite previously stored content, a failure mode termed catastrophic interference \cite{mccloskey1989catastrophic}. The interference is most severe precisely when successive contents are correlated and therefore recruit overlapping synaptic resources. Interleaving old content with new during learning limits the damage, since the shared substrate is constrained by both at once and converges to a solution compatible with each \cite{robins1995catastrophic}. The complementary learning systems account makes this the functional rationale for the two-structure organisation: a fast-learning hippocampal system encodes new episodes on the timescale of experience, and offline reactivation supplies to the slow-learning neocortical system the interleaved samples that it would otherwise have no way of obtaining, since the corresponding experiences are past \cite{mcclelland1995cls, kumaran2016what}. Under this hypothesis, replay is what makes it possible for a single cortical substrate to keep integrating new and potentially correlated content without erasing what it already holds. The way interleaving suppresses crosstalk between stored contents will be made explicit below in terms of the network dynamics.


[Now we need to talk about how this dialogue is considered to permit system consolidation and link it with memory transfer. if it is considered to do other things mention it while acknowledging we focus on the system consolidation/memory transfer stuff here. choose your world carefully.]

[ok now we need to detail mechanistically how this dialogue has been supposed to allow continuous learning. Here we have to acknowledge the various theories quickly (do some search) and to say that we will focus on the hypothesis which posit that throufh interleaved replay the brain find a way to minimize a phenomenon called interferences allowing for the same cortical substrat to continuously integrate new potentially correlated memory content without them to interfer with one another. we need to explain but not in too much details, the interleaving and removale of crosstalk/interferences will be illustrated later by the network dynamics.]

[Then we acknowledge that there have been proposed models of memory replay and system consolidation in bioplausible architectures/modesl in the litterature but that they tend to not explain which memory content would be prioritized for reactivation and how repeated reactivation relocates stored representations from one neuronal population to another, rather than \textbf{merely} training a downstream readout. Explain very well why it is very important to have model for both (the brain cannot just randomly pick up memories for replay it is not efficient as an exemple of justification). For this part you have to be very carefull do alot of researches and yes be diplomatic. maybe some models are more advanced that what I think. but anyway I think the core of my assertions is true.]

[ok now we start talking about our model ! first we explain that we are introducing a model based on the predictive coding framework which aim to offer an hypothesis on how engram reallocation and memory could happen in the brain using a bioinspired network using only local plasticity rules and dynamics.  
there is two big components to our model that are at its core. First we build on associative memory networks introduced in "Recurrent predictive coding models for associative memory employing covariance learning Mufeng Tang" which gives use the core properties for online engram reallocation by using the predictive coding framework in an associative memory task. 

We take two of these associative memory network , one a teacher that we can map to the hippocampus which has learned new memory content and a teacher which we can map to the neocortex toward which memory content will be reallocated. and we let the dialogue operate based on an interface error term classically borowed from predictive coding.

the second big component is that we depart significantly from classical predictive coding framework which aime solely at minimizing reconstruction error by changing the signe of the influence of the interface error term coordinating the dialogue between the student and the teacher.
Here we explain well that in a classical pc network this inteface error influence motivate to find a compromise between the influence of the different neuronal population. this classical influence tend to turn the system as a pattern completion system based on attractors/line attractors, sequence reconstructions ... bref it create a system that want to converge toward familiar activity pattern mostly, with maybe wandering due to noise. as we will show this behavior is not optimal for fast and exhaustive memory transfer , and more broadly for functional structured spontaneous activity as it means that the nervous system only try to revisit familiar activity pattern (acquired during awake activity) during sleep making it hard to explain an added computational property to what is already happening during awake.

ok no don't go into that much details, we will keep these details for discussion ! but yeah now you know it and you can put it in the discussion to be redacted later. 

bref so we talk about "this classical influence tend to turn the system as a pattern completion system based on attractors/line attractors, sequence reconstructions ... bref it create a system that want to converge toward familiar activity pattern mostly, with maybe wandering due to noise." and say then that by reversing the sign of the influence we strikingly (be more sober) the whole coupled system from a pattern completion system toward a system that orchestrate spontaneous activity to induce autonomous memory transfer via prioritized reactivation !  We call that sign reversal adversarial coupling.
Then quickly describe the behavior  of the system without going into to much details.
the teacher network stay near its memory manifold while via the adversarial coupling the student network bias the teacher network to explore part of the manifold that surprise it the most prioritizing the reactivation of poorly learned memory content (say that better). Then the predictive coding formulation of associative memory (tang paper) leads naturally to the learning of the expressed representions of the teacher during manifold exploration from the student via online updating of the recurrent synapses dure to reconstruction errors ! 
]

[then talk about the fact that we propose an architecture that uses this memory transfer dynamic as the building block for continuous learning via allowing the interleaved communication of 3 neuronal populations . one aking to the hippocampus learn quickly new memories, the second and the third aking to the neocortex (without precise mapping) integrate these new memory patterns while mitigating interferences even for strong correlations. precise that this architecture is not to be seen as a precise mapping or analogue of what happen in the brain but more of a demonstration of how we can use memory transfer dynamics to allow continuous learning via mitigation of interferences .]

[Then even if it will be deppen in the discussion say that the adversarial coupling is not out of nowhere, it is inspired by a recurrent pattern in neurobiology. corolarry discarges, subtractive top down signal (visual cortex, auditory to motor cortex put the citations) and finally recent observation on the influence of spontaneous activity in the prefrontal cortex on the activity of the hippocampus. but yeah it is mostly teasing will be mainly discussed in the discussion.]

%============================================================================
%  2. MODEL AND ANALYSIS   (Tang-style: model + theorems inline; proofs in Methods)
%============================================================================
\section{Model and analysis}

\subsection{A two-population predictive-coding memory}
% Teacher T (frozen, higher) above student S (plastic, lower); three error
% populations; interface error eps_TS = x_S - x_T; reversed precision pi_ST as
% the signed sleep/wake knob; free energies F_S, F_T. pi_TS > pi_S stated as an
% OPERATING REGIME, not a theorem assumption.

\subsection{Student inference and learning descend a free energy}
% Proposition 1 (stated here; proof in Materials and methods).
\begin{proposition}\label{prop:descent}
\textit{[Student state relaxation is exact gradient descent on $F_S$; the
zero-diagonal weight update is projected gradient descent on $F_S$.]}
\end{proposition}
% The zero-diagonal constraint is not ours to justify: descent under projection
% onto a linear subspace is the classical gradient-projection method
% \cite{levitinpolyak1966}, and the zero-diagonal recurrent weights (no autapses)
% are inherited from the covPCN substrate \cite{tang2023}. Equivalently, the
% diagonal is not a parameter, so learning is plain unconstrained descent on the
% off-diagonal weights; full algebra in Materials and methods.

\subsection{The student steers the teacher up the student's surprise}
% Theorem 1. Two structural facts, no operator algebra:
% (i) sign identity -- only the interface term of F_S contains x_T, so at ANY x_S
%     grad_{x_T} F_S = pi_TS (x_T - x_S) = -pi_TS eps_TS; the model's push on the
%     teacher is pi_ST eps_TS, so push = (|pi_ST|/pi_TS) grad_{x_T} F_S: teacher and
%     student feel the SAME interface energy through opposite-signed precisions.
% (ii) chain rule at the settled state -- x_S* is DEFINED by grad_{x_S} F_S = 0
%     (unique: the stationarity system is linear with positive-definite matrix),
%     so the chain-rule cross term vanishes identically and the partial gradient
%     at x_S* IS the full gradient of F_S^eq(x_T) = F_S(x_S*(x_T), x_T)
%     (= min_{x_S} F_S by strict convexity -- the variational reading, a remark
%     not a proof step). Classical envelope argument; only the chain rule used.
% No "teacher's surprise": F_T enters only through the self-pull.
\begin{theorem}\label{thm:prioritized}
\textit{[Fast student state, frozen student weights. Isolating the student's
contribution to the teacher's sleep dynamics --- the push $u=\pi_{ST}\varepsilon_{TS}$
--- the teacher performs exact gradient ascent on the student's settled variational
free energy $F_S^{\mathrm{eq}}(x_T)\equiv\min_{x_S}F_S(x_S,x_T)$:
$\tau_T\dot x_T|_{\mathrm{push}}=(|\pi_{ST}|/\pi_{TS})\,\nabla_{x_T}F_S^{\mathrm{eq}}(x_T)$
--- at every point, with no approximation, hence locally steepest ascent of the
student's surprise. Caveats: local ascent (not the global summit); the push rescales
but does not seed an empty direction (noise seeds it).]}
\end{theorem}

\subsection{The landscape: the novelty operator}
% Lemma 1: fast-student elimination -> N_S; annihilates learned directions.
% Proof is per-mode scalar (final_proofs.md §5): diagonalize S_S once, F_S
% separates into d one-variable parabolas, minimize each -> s_k* = (1-n_k)c_k.
% x_S* = (I - N_S) x_T: an attenuated copy (shrinkage, not a projector).
% Not needed for Theorem 1 -- it describes the landscape the theorem climbs.
% The per-mode teacher law (prioritization, c_k' ~ n_k c_k) follows directly from
% the modes -- no landscape needed; Cor 1 is where the two readings provably meet.
\begin{lemma}\label{lem:novelty}
\textit{[Fast-student steady state gives $x_S^*=(I-N_S)x_T$, hence
$\varepsilon_{TS}=-N_S x_T$, with $N_S=\pi_S S_S(\pi_{TS}I+\pi_S S_S)^{-1}$; $N_S$ is
symmetric PSD and $N_S u=0$ on learned directions.]}
\end{lemma}

\subsection{The surprise identity: the settled student's free energy is the novelty score}
% Corollary 1 (of Lemma 1). Bridges the two levels: novelty (weight-level operator N_S,
% eigenvalues n_k in [0,1)) is the spectrum of surprise (state-level scalar F_S).
% Direct computation: interface n^2 + self n(1-n) = n per direction.
\begin{corollary}\label{cor:surprise}
\textit{[At the fast-student equilibrium,
$F_S^{\mathrm{eq}}(x_T)=\tfrac{\pi_{TS}}{2}\,x_T^\top N_S\,x_T$: the landscape
Theorem~\ref{thm:prioritized} climbs is the novelty-weighted energy of the teacher's
state --- a quadratic that is flat on learned directions and curved on novel ones.]}
\end{corollary}
% Remark (prioritization -- the paper's namesake). Gradient ascent on a PSD quadratic
% decouples along the eigenvectors: in the student's frame, dc_k/dt =
% (|pi_ST|/tau_T) n_k c_k. Learned axes (n_k = 0) exactly frozen; novel axes amplified
% exponentially in proportion to novelty; among occupied axes the most novel dominates.
% Self-limiting: n_k <= (pi_S/pi_TS)||M_S u_k||^2, so the push on a direction dies as
% the student learns it -- the teacher climbs a landscape its own climbing flattens.
\begin{remark}\label{rem:prioritized}
\textit{[Axis by axis in the student's eigenframe, $\dot c_k=(|\pi_{ST}|/\tau_T)n_k
c_k$: learned directions frozen, novel directions amplified in proportion to their
novelty $n_k$; the push on each direction extinguishes as the student learns it.]}
\end{remark}

\subsection{Staying on the memory manifold}
% NO LEMMA HERE (cut 2026-07-20; archived in maths/additional_proofs_not_in_paper.md
% -- do not reintroduce). Operating choice + measurement instead: |pi_ST| kept
% well below pi_T (code default: half the stability threshold, pi_ST_safety=0.5);
% containment verified numerically (manifold_leakage, terminal_occupancy).
\textit{[We keep the reversed precision well below the teacher's self-precision
(half the measured stability threshold throughout); the teacher's self-pull then
dominates normal to the manifold, and we verify numerically that the teacher's
off-manifold occupancy remains negligible during sleep. This sets the
precision-weighting regime: strong enough to move the teacher along novel memory
directions, weak enough that the self-pull wins off the manifold.]}

\subsection{An exact free-energy saddle}
% Proposition 2: exact saddle iff pi_ST = -pi_TS; active-inference-REMINISCENT.
% One line from Theorem 1: self-pull = -grad F_T, push = (|pi_ST|/pi_TS) grad F_S^eq,
% so the single potential Phi = F_T - F_S^eq generates the teacher flow iff
% |pi_ST| = pi_TS (mobility matching).
\begin{proposition}\label{prop:saddle}
\textit{[The physical functional $\Phi=F_T-F_S$ generates the sleep dynamics (on
the teacher sphere, with the model's mobilities and projected $W_S$ ascent) iff
$\pi_{ST}=-\pi_{TS}$ --- by Theorem~\ref{thm:prioritized}, exactly when the teacher's
ascent rate on the student's surprise matches the precision $\pi_{TS}$ that defines it.
The teacher descends $\Phi$; the student, descending $F_S$
(Proposition~\ref{prop:descent}), ascends $\Phi$: a minimax on one potential. Shape
reminiscent of the active-inference agent--environment saddle, not identical.]}
\end{proposition}

\subsection{Scope: subspace, not named episodes}
\begin{remark}\label{rem:scope}
\textit{[Linearity $\Rightarrow$ transfer of a subspace, or effective rank, rather than named patterns. Discrete reactivation of individual memories requires an added nonlinearity; the generation of temporally ordered replay sequences lies outside the scope of the present model.]}
\end{remark}

%============================================================================
%  3. RESULTS   (simulations; declarative headings a la Tang)
%============================================================================
\section{Results}

\subsection{Prioritized transfer of unlearned directions}
% The primary result: restricted novelty spectrum falling; x_T alignment heatmap.
\figplaceholder{Fig~1. Novelty-spectrum ``staircase'' and $x_T$--manifold
alignment: the student acquires unlearned directions; ``reliable completion
across tested regimes'' (not global-convergence). Steps are a
timescale/plotting artifact.}

\subsection{The wake/sleep saddle picture}
% Bowl (predictable direction) vs hill (novel direction); learning flips hill->bowl.
\figplaceholder{Fig~2. Saddle cross-sections: a bowl along a predictable
direction vs.\ a hill along a novel one; the exact $\Phi=F_T-F_S$ picture at
$\pi_{ST}=-\pi_{TS}$.}

\subsection{Merging memories by interleaving}
% EMPIRICAL (no theorem): crosstalk sawtooth; interleaving merges teachers.
\figplaceholder{Fig~3. Interleaved rehearsal merges two teachers: the crosstalk
sawtooth decays to zero (union of memory subspaces), demonstrated in simulation.}

\subsection{Continual learning without catastrophic forgetting}
% EMPIRICAL: retention matrix grows vs buffer-only control collapsing to latest.
\figplaceholder{Fig~4. Continual learning: retention across a streamed sequence
vs.\ a buffer-only control that catastrophically forgets.}

%============================================================================
%  4. DISCUSSION
%============================================================================
\section{Discussion}
% - Free-energy interpretation: VFE mapping exact/derived; EFE mapping structural
%   (state honestly).
% - The saddle as ACTIVE-INFERENCE-REMINISCENT (shape, not identity).
% - Biological reading: hippocampus -> neocortex, sleep replay, wake/sleep = sign
%   of an interface precision.
% - Limitations: linear => subspace, not named episodes; nonlinear / soft-WTA
%   extension for episodic replay; completion is empirical, not proven.
% - Relation to Tang et al. and to our previous autonomous-retrieval paper.

%============================================================================
%  5. CONCLUSION
%============================================================================
\section{Conclusion}
% Short: transfer is the ownable primitive; it composes to union and continual
% learning; the linearity + PC framework make the mechanism exactly analyzable.

%============================================================================
%  MATERIALS AND METHODS   (full proofs + simulation details)
%============================================================================
\section*{Materials and methods}
% Full proofs of Proposition~\ref{prop:descent}, Theorem~\ref{thm:prioritized},
% Lemma~\ref{lem:novelty}, Corollary~\ref{cor:surprise},
% Proposition~\ref{prop:saddle}. (Source of record: ../maths/final_proofs.md —
% sober, self-contained, model + assumptions + all proofs, in the paper's order.
% Pedagogical long-form: detailed_proofs.md. Manifold containment: operating
% choice + numerical verification only — no lemma.)
% Adiabatic (fast-student) elimination; simulation conditions;
% parameters; numerical diagnostics (novelty_operator, circulation, spectral_gap,
% manifold_leakage). Future work: restricted-novelty (A_S) manifold analysis.

%============================================================================
%  SUPPORTING INFORMATION
%============================================================================
\section*{Supporting information}
% - No-noise slow-mixing version (annex).
% - Ablations; stop-grad / dendritic variant.
% - Correlation / sparsity / timescale sweeps.
% - Effective-rank / degeneracy discussion.

%============================================================================
%  REFERENCES
%============================================================================
\bibliographystyle{bibstyle}
\bibliography{references}

\end{document}
